\documentclass[11pt,a4paper]{article}

% ---------- Packages ----------
\usepackage[margin=2.2cm]{geometry}
\usepackage{booktabs}
\usepackage{longtable}
\usepackage{array}
\usepackage{tabularx}
\usepackage{listings}
\usepackage{xcolor}
\usepackage{hyperref}
\usepackage{enumitem}
\usepackage{titlesec}
\usepackage{fancyhdr}
\usepackage{parskip}
\usepackage{microtype}

% ---------- Code listing style ----------
\lstset{
    basicstyle=\ttfamily\small\color{black},
    breaklines=true,
    breakatwhitespace=false,
    frame=single,
    framesep=4pt,
    xleftmargin=4pt,
    xrightmargin=4pt,
    aboveskip=3pt,
    belowskip=3pt,
    showstringspaces=false,
    tabsize=4,
    columns=fullflexible,
    keepspaces=true,
    literate={-}{-}1
}

% Language definitions
\lstdefinestyle{python}{
    language=Python,
    keywordstyle=\color{blue}\bfseries,
    commentstyle=\color{green!50!black},
    stringstyle=\color{red!60!black}
}
\lstdefinestyle{javascript}{
    language=Java,
    morekeywords={const,let,require,module,exports,async,await,then,catch},
    keywordstyle=\color{blue}\bfseries,
    commentstyle=\color{green!50!black},
    stringstyle=\color{red!60!black}
}
\lstdefinestyle{dart}{
    language=Java,
    morekeywords={import,package,class,extends,implements,async,await,final,const,void,Future,String,bool,int,dynamic,override,super,State,StatefulWidget,StatelessWidget},
    keywordstyle=\color{blue}\bfseries,
    commentstyle=\color{green!50!black},
    stringstyle=\color{red!60!black}
}

% TypeScript/TSX style: blue keywords, green comments, red strings, black code
\lstdefinestyle{typescript}{
  language=Java,
  morekeywords={const,let,var,async,await,interface,type,export,default,import,from,as,extends,implements,return,if,else,switch,case,break,try,catch,throw,new,typeof,void,boolean,string,number,unknown,any,Promise,ReactNode,Readonly,enum,readonly,public,private,protected},
  keywordstyle=\color{blue}\bfseries,
  commentstyle=\color{green!50!black},
  stringstyle=\color{red!60!black}
}

% ---------- Header / Footer ----------
\pagestyle{fancy}
\fancyhf{}
\fancyhead[L]{\small NovaScan: AI-Driven DevSecOps Platform}
\fancyhead[R]{\small Compiled Source Code}
\fancyfoot[C]{\thepage}
\renewcommand{\headrulewidth}{0.4pt}

% ---------- Title formatting ----------
\titleformat{\section}{\large\bfseries}{\thesection}{1em}{}
\titleformat{\subsection}{\normalsize\bfseries}{\thesubsection}{1em}{}

\begin{document}

% ============================================================
% TITLE
% ============================================================
\begin{center}
    {\LARGE\bfseries NovaScan: AI-Driven DevSecOps Platform}\\[6pt]
    {\large Compiled Source Code Package for Software Copyright Registration}\\[2pt]
\end{center}

\vspace{1.2em}

% ============================================================
% 1. Software Copyright Identification & Metadata
% ============================================================
\section{Software Copyright Identification \& Metadata}

\begin{tabularx}{\textwidth}{@{}lX@{}}
\toprule
\textbf{TITLE OF WORK} & NovaScan: AI-Driven DevSecOps Platform \\
\midrule
\textbf{VERSION / DATE} & Repository snapshot / 14 November 2026 \\
\midrule
\textbf{CLASS OF WORK} & Computer Software / Computer Program\\
\midrule
\textbf{APPLICANT / OWNER} & Prakhar Shrivastava, Ayush Karanth, Rohan Bhangale, Suhas Arali, Rupali Kale
 \\
\midrule
\textbf{AUTHORS} & Prakhar Shrivastava, Ayush Karanth, Rohan Bhangale, Suhas Arali, Rupali Kale \\
\midrule
\textbf{PROGRAMMING LANGUAGES} & Python, JavaScript, TypeScript; Next.js 14 framework \\

\midrule
\textbf{INSTITUTION \& DEPARTMENT} & Vivekanand Education Society's Institute of Technology -- Information Technology Department \\
\bottomrule
\end{tabularx}

\vspace{1.5em}

% ============================================================
% 2. Statement of Third-Party Material
% ============================================================
\section{Statement of Third-Party Material \& Open Source Dependencies}

{\small
\begin{longtable}{@{}p{3.1cm}p{3.5cm}p{2.1cm}p{1.5cm}p{5.1cm}@{}}
\toprule
\textbf{Component} & \textbf{Source / Repository} & \textbf{Licence} & \textbf{Use} & \textbf{Scope / Attribution} \\
\midrule
\endfirsthead
\toprule
\textbf{Component} & \textbf{Source / Repository} & \textbf{Licence} & \textbf{Use} & \textbf{Scope / Attribution} \\
\midrule
\endhead
\bottomrule
\endfoot
FastAPI & PyPI / fastapi & MIT & Direct dependency & Backend REST API framework. \\
\midrule
Uvicorn & PyPI / uvicorn & BSD-3-Clause & Direct dependency & ASGI server. \\
\midrule
Pydantic, pydantic-settings & PyPI & MIT & Direct dependencies & Request validation and configuration. \\
\midrule
Motor, PyMongo & PyPI & Apache-2.0 & Direct dependencies & MongoDB driver and BSON support. \\
\midrule
Redis Python client & PyPI / redis & MIT & Direct dependency & Redis client library. \\
\midrule
python-jose & PyPI & MIT & Direct dependency & JWT handling. \\
\midrule
Passlib, bcrypt & PyPI & BSD-3-Clause / Apache-2.0 & Direct dependencies & Password hashing utilities. \\
\midrule
python-multipart & PyPI & Apache-2.0 & Direct dependency & Multipart form and upload parsing. \\
\midrule
pypdf & PyPI & BSD-3-Clause & Direct dependency & PDF text extraction. \\
\midrule
google-auth & PyPI & Apache-2.0 & Direct dependency & Google OAuth credential support. \\
\midrule
google-genai, google-generativeai & PyPI & Apache-2.0 & Direct dependencies & Google AI client libraries. \\
\midrule
slack-sdk & PyPI & MIT & Direct dependency & Slack API client. \\
\midrule
httpx & PyPI & BSD-3-Clause & Direct dependency & HTTP client for external APIs. \\
\midrule
python-dotenv & PyPI & BSD-3-Clause & Direct dependency & Environment-file loading. \\
\midrule
jira & PyPI & BSD-2-Clause & Direct dependency & Jira client library. \\
\midrule
resend & PyPI & MIT & Direct dependency & Resend email API client. \\
\midrule
email-validator & PyPI & Unlicense & Direct dependency & Email address validation. \\
\midrule
tiktoken & PyPI & MIT & Direct dependency & Token counting utility. \\
\midrule
tenacity & PyPI & Apache-2.0 & Direct dependency & Retry utilities. \\
\midrule
Next.js & npm / next & MIT & Direct dependency & Frontend application framework (14.2.35). \\
\midrule
React, React DOM & npm & MIT & Direct dependencies & UI rendering and component runtime. \\
\midrule
Tailwind CSS & npm / tailwindcss & MIT & Direct dependency & Utility CSS framework. \\
\midrule
Base UI & npm / @base-ui/react & MIT & Direct dependency & UI primitives. \\
\midrule
Radix UI React packages & npm / @radix-ui/react-* & MIT & Direct dependencies & Accessible UI primitives used by the frontend. \\
\midrule
@react-oauth/google & npm & MIT & Direct dependency & Google sign-in UI integration. \\
\midrule
TanStack React Table & npm / @tanstack/react-table & MIT & Direct dependency & Table state and rendering utilities. \\
\midrule
Axios & npm & MIT & Direct dependency & Frontend HTTP client. \\
\midrule
class-variance-authority & npm & Apache-2.0 & Direct dependency & Component variant utility. \\
\midrule
clsx, cn & npm & MIT & Direct dependencies & Class-name utilities. \\
\midrule
Framer Motion & npm & MIT & Direct dependency & UI animation library. \\
\midrule
jwt-decode & npm & MIT & Direct dependency & JWT payload decoding in the frontend. \\
\midrule
Lucide React & npm / lucide-react & ISC & Direct dependency & UI icon set. \\
\midrule
next-themes & npm & MIT & Direct dependency & Theme preference support. \\
\midrule
Recharts & npm & MIT & Direct dependency & Data visualization components. \\
\midrule
tailwind-merge, tailwindcss-animate & npm & MIT & Direct dependencies & Tailwind class and animation utilities. \\
\midrule
TypeScript & npm / typescript & Apache-2.0 & Dev dependency & Frontend language tooling. \\
\midrule
ESLint, eslint-config-next, PostCSS & npm & MIT & Dev dependencies & Linting and CSS build tooling. \\
\midrule
@types/node, @types/react, @types/react-dom & npm & MIT & Dev dependencies & TypeScript declaration packages. \\
\midrule
GitHub App and REST APIs & GitHub & GitHub terms & External service & Repository installation, pull request metadata, and webhook delivery; API service, not bundled source code. \\
\midrule
Google OAuth, Calendar, Drive, Meet, and Gemini APIs & Google & Google API terms & External services & OAuth and external Google platform APIs; access depends on configured scopes and Google project status. \\
\midrule
Slack Web API & Slack & Slack developer terms & External service & Workspace messaging integration. \\
\midrule
Atlassian Jira Cloud REST API & Atlassian & Atlassian terms & External service & Jira issue tracking integration. \\
\midrule
Notion API & Notion & Notion API terms & External service & Audit-report and workspace content integration. \\
\midrule
MongoDB service & MongoDB & MongoDB terms & External service & Application database hosting/service, where configured. \\
\bottomrule
\end{longtable}

\vspace{0.6em}
The dependency rows above summarize direct packages declared in \texttt{backend/requirements.txt} and \texttt{frontend/package.json}; npm resolutions are recorded in frontend/package-lock.json; Python dependency versions are declared in backend/requirements.txt, with unpinned entries resolved by the installer. Licence identifiers are shown as declared by upstream project metadata and should be confirmed against the exact distributed versions before filing.

Third-party platform integrations are called as external services. Their APIs, hosted services, product names, and terms remain owned by their respective providers; this compilation claims no ownership of those materials. The application source code contains integration logic and configuration references, not copies of provider API implementations.

\vspace{0.8em}
\noindent\textit{Declaration:} The applicants claim copyright only in original source code authored by them. Third-party software dependencies remain subject to their respective licences and notices. No ownership claim is made over third-party libraries, APIs, hosted services, or pretrained models.

\vspace{1.5em}

% ============================================================
% 3. Repository Source Code Listings
% ============================================================
\section{Selected Source Code Excerpts}

The following excerpts are taken from ten repository files in the 14 November 2026 snapshot. They show representative original implementation portions; omitted portions of each file are not reproduced. Unicode symbols in comments or display strings are rendered as ASCII equivalents for LaTeX portability.

\subsection{FILE 1: Backend application entry point\\(backend/main.py)}

\textit{Selected excerpt. Application lifespan, database connections, CORS configuration, and API route registration.}

\begin{lstlisting}[style=python]
@asynccontextmanager
async def lifespan(app: FastAPI):
    # Startup: Connect to DBs
    await connect_to_mongo()
    await connect_to_redis()
    yield
    # Shutdown: Close connections
    await close_mongo_connection()
    await close_redis_connection()
# --- INITIALIZE APP ONCE ---
app = FastAPI(
    title=settings.PROJECT_NAME,
    lifespan=lifespan,
    openapi_url=f"{settings.API_V1_STR}/openapi.json"
)
origins = [
    "http://localhost:3000",
    "http://127.0.0.1:3000",
    "https://nova-scan.vercel.app",
]
if getattr(settings, "FRONTEND_URL", None) and settings.FRONTEND_URL not in origins:
    origins.append(settings.FRONTEND_URL)
app.add_middleware(
    CORSMiddleware,
    allow_origins=origins,
    allow_credentials=True,
    allow_methods=["*"],
    allow_headers=["*"],
)
# --- WIRE UP THE ROUTES ---
app.include_router(api_router, prefix=settings.API_V1_STR)
\end{lstlisting}

\subsection{FILE 2: Authentication and security helpers\\(backend/app/core/security.py)}

\textit{Selected source. Password hashing and complete JWT access-token creation.}

\begin{lstlisting}[style=python]
from datetime import datetime, timedelta
from typing import Optional
from jose import jwt
from passlib.context import CryptContext
from app.core.config import settings
# Password Hashing Setup
pwd_context = CryptContext(schemes=["bcrypt"], deprecated="auto")
def verify_password(plain_password: str, hashed_password: str) -> bool:
    return pwd_context.verify(plain_password, hashed_password)
def get_password_hash(password: str) -> str:
    return pwd_context.hash(password)
# JWT Token Logic
def create_access_token(data: dict, expires_delta: Optional[timedelta] = None):
    to_encode = data.copy()
    if expires_delta:
        expire = datetime.utcnow() + expires_delta
    else:
        expire = datetime.utcnow() + timedelta(minutes=15)
    to_encode.update({"exp": expire})
    encoded_jwt = jwt.encode(to_encode, settings.SECRET_KEY, algorithm="HS256")
    return encoded_jwt
\end{lstlisting}

\subsection{FILE 3: Authentication endpoints\\(backend/app/api/v1/endpoints/auth.py)}

\textit{Selected excerpt. Privacy-preserving password reset request, hashed reset tokens, expiry, and password update.}

\begin{lstlisting}[style=python]
@router.post("/forgot-password")
async def forgot_password(payload: ForgotPasswordRequest):
    """Always return the same response so account existence is not disclosed."""
    generic_message = "If an account exists for that email, we've sent a password reset link."
    db = get_database()
    user = await db["users"].find_one({"email": payload.email})
    if not user:
        return {"message": generic_message}
    raw_token = secrets.token_urlsafe(32)
    token_hash = hashlib.sha256(raw_token.encode("utf-8")).hexdigest()
    expires_at = datetime.utcnow() + timedelta(minutes=30)
    await db["users"].update_one(
        {"_id": user["_id"]},
        {"$set": {"password_reset_token_hash": token_hash, "password_reset_expires_at": expires_at}},
    )
    reset_url = f"{settings.FRONTEND_URL.rstrip('/')}/auth/reset-password?{urlencode({'token': raw_token})}"
    sent = await send_password_reset_email(payload.email, reset_url)
    if not sent:
        logger.error("Password reset email could not be delivered for an account")
    return {"message": generic_message}
@router.post("/reset-password")
async def reset_password(payload: ResetPasswordRequest):
    token_hash = hashlib.sha256(payload.token.encode("utf-8")).hexdigest()
    db = get_database()
    user = await db["users"].find_one({
        "password_reset_token_hash": token_hash,
        "password_reset_expires_at": {"$gt": datetime.utcnow()},
    })
    if not user:
        raise HTTPException(status_code=400, detail="This reset link is invalid or expired. Request a new one.")
    result = await db["users"].update_one(
        {
            "_id": user["_id"],
            "password_reset_token_hash": token_hash,
            "password_reset_expires_at": {"$gt": datetime.utcnow()},
        },
        {
            "$set": {"hashed_password": security.get_password_hash(payload.new_password)},
            "$unset": {"password_reset_token_hash": "", "password_reset_expires_at": ""},
        },
    )
    if result.modified_count != 1:
        raise HTTPException(status_code=400, detail="This reset link is invalid or expired. Request a new one.")
    return {"message": "Your password has been reset. You can now sign in."}
\end{lstlisting}

\textit{Additional excerpt: Authenticated password change}

\begin{lstlisting}[style=python]
async def change_my_password(
    payload: PasswordChangeRequest,
    current_user: User = Depends(get_current_user),
):
    db = get_database()
    user = await db["users"].find_one({"email": current_user.email})
    if not user:
        raise HTTPException(status_code=404, detail="User account not found")
    existing_hash = user.get("hashed_password")
    if existing_hash:
        if not payload.current_password:
            raise HTTPException(status_code=400, detail="Enter your current password")
        if not security.verify_password(payload.current_password, existing_hash):
            raise HTTPException(status_code=400, detail="Current password is incorrect")
        if security.verify_password(payload.new_password, existing_hash):
            raise HTTPException(status_code=400, detail="Choose a password different from your current one")
    await db["users"].update_one(
        {"_id": user["_id"]},
        {"$set": {"hashed_password": security.get_password_hash(payload.new_password)}},
    )
    return {"message": "Password updated successfully"}
\end{lstlisting}

\subsection{FILE 4: GitHub webhook project scoping\\(backend/app/api/v1/endpoints/webhooks.py)}

\textit{Selected excerpt. Resolves webhook credentials from the matching project and verifies the repository identity.}

\begin{lstlisting}[style=python]
async def get_org_credentials(org_id: str, repo_full_name: str = None, project_id: str = None):
    """Fetch integration secrets from the matching project only; never org-wide."""
    try:
        db = get_database()
        project = None
        if project_id:
            query = {"_id": project_id, "organization_id": org_id}
            if ObjectId.is_valid(project_id):
                query = {"_id": ObjectId(project_id), "organization_id": org_id}
            project = await db["projects"].find_one(query)
        elif repo_full_name:
            owner, _, repo_name = repo_full_name.partition("/")
            project = await db["projects"].find_one({
                "organization_id": org_id,
                "repo_owner": {"$regex": f"^{re.escape(owner)}$", "$options": "i"},
                "repo_name": {"$regex": f"^{re.escape(repo_name)}$", "$options": "i"},
            })
        if not project:
            logger.warning("No project found for integration event (%s)", repo_full_name or project_id)
            return None
        if repo_full_name:
            expected_repo = f"{project.get('repo_owner', '')}/{project.get('repo_name', '')}"
            if expected_repo.casefold() != repo_full_name.casefold():
                logger.warning("Ignoring repo event that does not match project '%s'", project.get("name"))
                return None
        settings_data = project.get("settings", {})
        oauth_credentials = await oauth_token_manager.credentials_for_project(org_id, str(project["_id"]))
        # GitHub API access must come from the repo-scoped GitHub App token.
        github_token = oauth_credentials.get("github_token")
        return {
            "project_id": str(project["_id"]),
            "github_token": github_token,
\end{lstlisting}

\subsection{FILE 5: Risk analysis service\\(backend/app/services/intelligence/risk\_analyzer.py)}

\textit{Selected excerpt. Builds event-specific analysis instructions, invokes the shared model service, and handles failures.}

\begin{lstlisting}[style=python]
from app.services.intelligence.llm_engine import generate_insight
async def analyze_risk(event: dict) -> dict:
    """
    Analyzes an incoming event (GitHub PR, Jira Ticket, Slack Msg) using Gemini
    to determine if it poses a risk to the project.
    """
    # 1. Select the Best Persona based on Source
    if event["source"] == "github":
        persona = "You are a Senior Staff Engineer. Look for breaking API changes, security flaws, risky dependencies, or lack of tests."
    elif event["source"] == "jira":
        persona = "You are a Technical Product Manager. Look for vague requirements, scope creep, sudden priority spikes, or unrealistic deadlines."
    elif event["source"] == "slack":
        persona = "You are an Agile Delivery Lead. Look for signs of team blocking, miscommunication, delays, or frustration."
    else:
        persona = "You are a Project Risk Analyst."
    # 2. Construct the System Prompt
    system_instruction = f"""
    {persona}
    Your task is to analyze the following project event and determine if it requires attention.
    Output strictly in this JSON format:
    {{
        "is_risky": boolean,
        "risk_score": integer (1-10),
        "reason": "A concise, one-sentence explanation of the risk.",
        "action_needed": "A short, actionable recommendation for the team."
    }}
    Rules:
    - If risk_score < 4, set is_risky to false.
    - If code changes involve database schemas, authentication, or payments, risk_score must be at least 8.
    - Be concise.
    """
    # 3. Construct the User Content (The Data)
    user_content = f"""
    Event Source: {event['source']}
    Action: {event['action']}
    Actor: {event['actor']}
    Summary: {event['summary']}
    Full Content/Diff/Message:
    {event['content']}
    """
    # 4. Call the Shared LLM Engine
    try:
        # We use the generic wrapper we built in llm_engine.py
        analysis = await generate_insight(system_instruction, user_content)
        return analysis
    except Exception as e:
        print(f"[WARN] Risk Analyzer Error: {e}")
        # Fail safe (Default to safe so we don't spam alerts on errors)
        return {
            "is_risky": False,
            "risk_score": 0,
            "reason": "AI Analysis temporarily unavailable",
            "action_needed": "Manual review recommended"
        }
\end{lstlisting}

\subsection{FILE 6: Repository-scoped GitHub token\\(backend/app/services/integration/github\_app.py)}

\textit{Selected excerpt. Requests a short-lived GitHub App token scoped to one repository and read permissions.}

\begin{lstlisting}[style=python]
async def create_repository_token(installation_id: str, repository_id: str) -> str:
    """Mint an ephemeral token that can access exactly the supplied repository."""
    headers = {
        "Authorization": f"Bearer {_app_jwt()}",
        "Accept": "application/vnd.github+json",
        "X-GitHub-Api-Version": "2022-11-28",
    }
    async with httpx.AsyncClient(timeout=20.0) as client:
        response = await client.post(
            f"https://api.github.com/app/installations/{installation_id}/access_tokens",
            headers=headers,
            json={
                "repository_ids": [int(repository_id)],
                "permissions": {"contents": "read", "pull_requests": "read"},
            },
        )
    if response.status_code != 201:
        raise RuntimeError(f"Could not create repository-scoped GitHub token ({response.status_code})")
    token = response.json().get("token")
    if not token:
        raise RuntimeError("GitHub did not return an installation token")
    return token
\end{lstlisting}

\textit{Additional excerpt: Builds the short-lived signed GitHub App JWT}

\begin{lstlisting}[style=python]
def _app_jwt() -> str:
    encoded_key = getattr(settings, "GITHUB_APP_PRIVATE_KEY_BASE64", None)
    raw_key = getattr(settings, "GITHUB_APP_PRIVATE_KEY", None)
    if not (settings.GITHUB_APP_ID and (encoded_key or raw_key)):
        raise RuntimeError("GitHub App ID and private key are not configured")
    if encoded_key:
        try:
            private_key = base64.b64decode(re.sub(r"\s+", "", str(encoded_key)), validate=True).decode("utf-8-sig")
        except (binascii.Error, UnicodeDecodeError, ValueError) as exc:
            raise RuntimeError("GITHUB_APP_PRIVATE_KEY_BASE64 must be Base64-encoded UTF-8 PEM file contents") from exc
    else:
        private_key = str(raw_key).strip()
    if len(private_key) >= 2 and private_key[0] == private_key[-1] and private_key[0] in ("\"", "'"):
        private_key = private_key[1:-1].strip()
    private_key = (
        private_key
        .replace("\\r\\n", "\n")
        .replace("\\n", "\n")
        .replace("\r\n", "\n")
        .replace("\r", "\n")
    )
    # Render and other environment managers may preserve a pasted PEM as one
    # long line. Canonicalize only a recognizable, unencrypted PEM block.
    pem = re.search(r"-----BEGIN ((?:RSA |EC )?PRIVATE KEY)-----([\s\S]*?)-----END \1-----", private_key)
    if pem:
        body = re.sub(r"\s+", "", pem.group(2))
        if body and re.fullmatch(r"[A-Za-z0-9+/=]+", body):
            private_key = (
                f"-----BEGIN {pem.group(1)}-----\n"
                + "\n".join(body[index:index + 64] for index in range(0, len(body), 64))
                + f"\n-----END {pem.group(1)}-----\n"
            )
    now = datetime.now(timezone.utc)
    try:
        return jwt.encode(
            {"iat": int((now - timedelta(seconds=30)).timestamp()),
             "exp": int((now + timedelta(minutes=8)).timestamp()),
             "iss": settings.GITHUB_APP_ID},
            private_key,
            algorithm="RS256",
        )
    except Exception as exc:
        raise RuntimeError(
            "GitHub App private key is not valid PEM. In Render, replace GITHUB_APP_PRIVATE_KEY with the complete, unencrypted .pem file contents, including both BEGIN/END lines."
        ) from exc
\end{lstlisting}

\subsection{FILE 7: Google Calendar OAuth state validation\\(backend/app/api/v1/endpoints/standup.py)}

\textit{Selected excerpt. Signs and validates OAuth state and restricts callback and frontend return URLs.}

\begin{lstlisting}[style=python]
def _google_state_encode(user_id: str, org_id: str, team_id: str, redirect_uri: str, frontend_origin: str, include_drive: bool) -> str:
    return jwt.encode({"sub": str(user_id), "org_id": str(org_id), "team_id": team_id, "redirect_uri": redirect_uri, "frontend_origin": frontend_origin, "include_drive": include_drive, "exp": datetime.utcnow() + timedelta(minutes=15)}, settings.SECRET_KEY, algorithm="HS256")
async def _google_state_decode(state: str):
    try:
        claims = jwt.decode(state, settings.SECRET_KEY, algorithms=["HS256"])
    except JWTError as exc:
        raise HTTPException(status_code=400, detail="Google Calendar authorization expired. Start the connection again.") from exc
    if not claims.get("sub") or not claims.get("org_id") or not claims.get("team_id") or not claims.get("redirect_uri") or not claims.get("frontend_origin"):
        raise HTTPException(status_code=400, detail="Invalid Google Calendar authorization state")
    redirect_uri = claims["redirect_uri"]
    allowed_redirects = {settings.google_calendar_redirect}
    parsed_redirect = urlparse(redirect_uri)
    if parsed_redirect.scheme == "http" and parsed_redirect.hostname in ("localhost", "127.0.0.1", "::1") and parsed_redirect.path == "/api/v1/standup/google/callback" and not parsed_redirect.query and not parsed_redirect.fragment:
        allowed_redirects.add(redirect_uri)
    if redirect_uri not in allowed_redirects:
        raise HTTPException(status_code=400, detail="Google Calendar callback URI is not allowed")
    frontend_origin = claims["frontend_origin"].rstrip("/")
    parsed_frontend = urlparse(frontend_origin)
    if frontend_origin != settings.FRONTEND_URL.rstrip("/") and not (parsed_frontend.scheme == "http" and parsed_frontend.hostname in ("localhost", "127.0.0.1") and parsed_frontend.port == 3000 and not parsed_frontend.path and not parsed_frontend.query and not parsed_frontend.fragment):
        raise HTTPException(status_code=400, detail="Google Calendar return URL is not allowed")
    return claims
\end{lstlisting}

\textit{Additional excerpt: Validates the callback and returns the user to standups}

\begin{lstlisting}[style=python]
async def google_calendar_callback(code: Optional[str] = None, state: Optional[str] = None, error: Optional[str] = None):
    frontend = settings.FRONTEND_URL.rstrip("/") + "/dashboard/standup"
    if error or not code or not state:
        return RedirectResponse(frontend + "?" + urlencode({"google": "error", "message": error or "Authorization was cancelled"}))
    try:
        claims = await _google_state_decode(state)
        frontend = claims["frontend_origin"].rstrip("/") + "/dashboard/standup"
        db = get_database()
        user = await db["users"].find_one({"_id": {"$in": [claims["sub"], ObjectId(claims["sub"])]} if ObjectId.is_valid(claims["sub"]) else claims["sub"]})
        if not user or str(user.get("current_org_id")) != str(claims["org_id"]):
            raise HTTPException(status_code=403, detail="Google Calendar account no longer belongs to this organization")
        team = await _find_team_for_standup(db, claims["org_id"], claims["team_id"])
        if not team:
            raise HTTPException(status_code=404, detail="Team not found")
        callback_user = User(**{**user, "_id": str(user.get("_id"))})
        if not _team_manage_allowed(team, callback_user):
            raise HTTPException(status_code=403, detail="You no longer have permission to connect this team's calendar")
        token_data = await google_calendar.exchange_code(code, claims["redirect_uri"])
        google_email = await google_calendar.google_account_email(token_data.get("access_token", ""))
        await google_calendar.save_connection(claims["org_id"], claims["sub"], token_data)
        await db["google_calendar_connections"].update_one(
            {"org_id": claims["org_id"], "user_id": claims["sub"]},
            {"$set": {"email": google_email, "app_user_email": user.get("email"), "full_name": user.get("full_name") or user.get("name")}},
        )
        return RedirectResponse(frontend + "?" + urlencode({"google": "drive_connected" if claims.get("include_drive") else "connected", "team_id": claims["team_id"]}))
    except HTTPException as exc:
        logger.warning("Google Calendar OAuth callback rejected: %s", exc.detail)
        return RedirectResponse(frontend + "?" + urlencode({"google": "error", "message": exc.detail}))
    except (ValueError, httpx.HTTPError) as exc:
        logger.warning("Google Calendar OAuth callback failed: %s", exc)
        return RedirectResponse(frontend + "?google=error&message=Google+Calendar+connection+failed")
\end{lstlisting}

\subsection{FILE 8: Standup transcript analysis\\(backend/app/services/intelligence/standup\_agent.py)}

\textit{Selected excerpt. Provides team context to the model and extracts tasks, events, blockers, and a summary.}

\begin{lstlisting}[style=python]
import google.generativeai as genai
import os
import json
import logging
from typing import List, Dict
from app.models.user import User
# Configure Logging
logger = logging.getLogger(__name__)
# Load API Key
GEMINI_API_KEY = os.getenv("GEMINI_API_KEY")
if GEMINI_API_KEY:
    genai.configure(api_key=GEMINI_API_KEY)
MODEL_NAME = "gemini-2.5-flash-lite"
async def analyze_standup(transcript: str, members: List[User], project_name: str) -> Dict:
    """
    Analyzes a standup transcript to identify speakers, action items, events, and blockers.
    """
    if not GEMINI_API_KEY:
        logger.error("[ERROR] GEMINI_API_KEY is missing!")
        return {"error": "AI Service Unavailable"}
    # Create a mapping string for the LLM
    member_list_str = "\n".join([f"- {u.full_name} (ID: {u.id}, Email: {u.email})" for u in members])
    prompt = f"""
    You are an AI Project Manager. Analyze this Daily Standup Transcript for project "{project_name}".
    Match speakers in the text to these Team Members:
    {member_list_str}
    TRANSCRIPT:
    {transcript}
    ---------------------------------------------------
    Task:
    1. Identify who is speaking. Map names in text to the provided Team Members.
    2. Extract 'tasks' (Action Items). Assign to the correct User ID.
    3. Extract 'events' (Deadlines, Meetings). accurate dates/times if mentioned.
    4. Extract 'blockers' (Who is stuck and why).
    5. Generate a 'summary' (Markdown format) for Slack.
    Output JSON ONLY:
    {{
        "tasks": [
            {{ "title": "...", "description": "...", "assignee_id": "<User ID from list>", "assignee_name": "..." }}
        ],
        "events": [
            {{ "title": "...", "date": "YYYY-MM-DD or 'Today'", "time": "HH:MM", "description": "..." }}
        ],
        "blockers": [
            {{ "user_id": "<User ID>", "user_name": "...", "issue": "..." }}
        ],
        "summary": "### [START] Daily Standup Summary - {project_name}\\n\\n**Highlights**... "
    }}
    """
    try:
        model = genai.GenerativeModel(MODEL_NAME)
        response = model.generate_content(prompt)
        text_response = response.text.strip()
        if text_response.startswith("```"):
            text_response = text_response.strip("`").replace("json", "").strip()
        return json.loads(text_response)
    except Exception as e:
        logger.error(f"[ERROR] Standup Analysis Failed: {e}")
        return {"error": str(e)}
\end{lstlisting}

\subsection{FILE 9: Notion audit report construction\\(backend/app/services/delivery/notion\_bot.py)}

\textit{Selected excerpt. Builds report blocks and creates a page under the selected parent page or database.}

\begin{lstlisting}[style=python]
        children = [{
            "object": "block", "type": "callout", "callout": {
                "rich_text": self._rich_text(f"NovaScan risk audit | {self._get_risk_label(risk_score)} | Score {risk_score}/10"),
                "icon": {"type": "emoji", "emoji": "\U0001F4CB"},
                "color": "red_background" if risk_score >= 7 else "orange_background",
            }
        }, heading("Executive risk assessment"), *paragraphs(analysis.get("reason") or analysis.get("summary")), heading("Recommended actions")]
        for item in recommendations:
            children.append({"object": "block", "type": "bulleted_list_item", "bulleted_list_item": {"rich_text": self._rich_text(item)}})
        for key, label in (("impact", "Potential impact"), ("affected_areas", "Affected areas"), ("evidence", "Evidence"), ("additional_recommendations", "Additional recommendations")):
            value = analysis.get(key)
            if value:
                children.append(heading(label, 3))
                for item in value if isinstance(value, list) else [value]:
                    children.append({"object": "block", "type": "bulleted_list_item", "bulleted_list_item": {"rich_text": self._rich_text(item)}})
        children.extend([
            heading("Event details"),
            *paragraphs(f"Repository: {context.get('repo') or 'Unknown'}"),
            *paragraphs(f"Pull request: #{context.get('pr_id') or 'N/A'} | {context.get('title') or event_summary}"),
            *paragraphs(f"Author: {context.get('author') or 'Unknown'}"),
            *paragraphs(f"Detected at (UTC): {created_at}"),
        ])
        if isinstance(pr_url, str) and pr_url.startswith("http"):
            children.append({"object": "block", "type": "paragraph", "paragraph": {"rich_text": [{"type": "text", "text": {"content": "Open source pull request", "link": {"url": pr_url}}}]}})
        if parent_page_id:
            payload = {
                "parent": {"page_id": parent_page_id},
                "properties": {"title": {"title": [{"type": "text", "text": {"content": title}}]}},
                "children": children[:100],
            }
\end{lstlisting}

\textit{Additional excerpt: Creates the parented report and handles Notion API responses}

\begin{lstlisting}[style=python]
        else:
            payload = {
                "parent": {"database_id": database_id},
                "properties": {
                    "Name": {"title": [{"text": {"content": title}}]},
                    "Risk Level": {"select": {"name": self._get_risk_label(risk_score)}},
                    "Risk Score": {"number": risk_score},
                    "Status": ({"select": {"name": "In progress"}} if database_status_type == "select" else {"status": {"name": "In Progress"}}),
                    "Source PR": {"url": pr_url if isinstance(pr_url, str) and pr_url.startswith("http") else "https://github.com"},
                    "Detected At": {"date": {"start": created_at}},
                },
                "children": children[:100],
            }
        try:
            async with httpx.AsyncClient(timeout=25.0) as client:
                response = await client.post(f"{self.base_url}/pages", headers=self._get_headers(notion_token), json=payload)
                if response.status_code not in (200, 201):
                    logger.error("Notion audit page creation failed (%s): %s", response.status_code, response.text[:500])
                    return None
                return response.json()
        except Exception as exc:
            logger.error("Notion audit page creation failed: %s", exc)
            return None
\end{lstlisting}

\subsection{FILE 10: Global timezone display helpers\\(frontend/app/dashboard/standup/page.tsx)}

\textit{Selected excerpt. Formats local times and validates IANA timezone identifiers in the standup UI.}

\begin{lstlisting}[style=typescript]
function timezoneOffset(timeZone: string) {
  try { return new Intl.DateTimeFormat(undefined, { timeZone, timeZoneName: "shortOffset" }).formatToParts(new Date()).find((part) => part.type === "timeZoneName")?.value || "GMT"; }
  catch { return "GMT"; }
}
function formattedLocalTime(timeZone: string) {
  try { return new Intl.DateTimeFormat(undefined, { timeZone, hour: "numeric", minute: "2-digit", timeZoneName: "short" }).format(new Date()); }
  catch { return ""; }
}
function isValidTimeZone(timeZone: string) {
  try { new Intl.DateTimeFormat(undefined, { timeZone }).format(); return true; }
  catch { return false; }
}
\end{lstlisting}

\newpage

\end{document}
